%%=============================================================================
%% Samenvatting
%%=============================================================================

% TODO: De "abstract" of samenvatting is een kernachtige (~ 1 blz. voor een
% thesis) synthese van het document.
%
% Een goede abstract biedt een kernachtig antwoord op volgende vragen:
%
% 1. Waarover gaat de bachelorproef?
% 2. Waarom heb je er over geschreven?
% 3. Hoe heb je het onderzoek uitgevoerd?
% 4. Wat waren de resultaten? Wat blijkt uit je onderzoek?
% 5. Wat betekenen je resultaten? Wat is de relevantie voor het werkveld?
%
% Daarom bestaat een abstract uit volgende componenten:
%
% - inleiding + kaderen thema
% - probleemstelling
% - (centrale) onderzoeksvraag
% - onderzoeksdoelstelling
% - methodologie
% - resultaten (beperk tot de belangrijkste, relevant voor de onderzoeksvraag)
% - conclusies, aanbevelingen, beperkingen
%
% LET OP! Een samenvatting is GEEN voorwoord!

%%---------- Nederlandse samenvatting -----------------------------------------
%
% TODO: Als je je bachelorproef in het Engels schrijft, moet je eerst een
% Nederlandse samenvatting invoegen. Haal daarvoor onderstaande code uit
% commentaar.
% Wie zijn bachelorproef in het Nederlands schrijft, kan dit negeren, de inhoud
% wordt niet in het document ingevoegd.

\IfLanguageName{english}{%
\selectlanguage{dutch}
\chapter*{Samenvatting}
\lipsum[1-4]
\selectlanguage{english}
}{}

%%---------- Samenvatting -----------------------------------------------------
% De samenvatting in de hoofdtaal van het document

\chapter*{\IfLanguageName{dutch}{Samenvatting}{Abstract}}

Deze bachelorthesis toont het onderzoek naar hoe een geografisch informatiesysteem (GIS)-gebaseerde tool kan bijdragen aan de energieplanning voor de Einstein Telescope (ET), een geplande geavanceerde zwaartekrachtgolfdetector waarvoor de Euregio Maas-Rijn (EMR) een kandidaat locatie is. De ET zal een aanzienlijk energieverbruik hebben, vergelijkbaar met dat van UZ Leuven of CERN. Hierbij stelt de ET ook hoge eisen aan de energievoorziening op het vlak van betrouwbaarheid, trillingsvrijheid en duurzaamheid. Daarnaast moet ook rekening worden gehouden met het lokale maatschappelijke draagvlak en wensen ze ook in te zetten op innovatieve hernieuwbare energiebronnen zoals agrophotovoltaïcs (AgriPV).\par

De grensoverschrijdende regio, ontbreken van een exacte locatie en ecologische beperkingen (bv. Natura2000-gebieden) maken traditionele energieanalyses tijdrovend en foutgevoelig. Deze thesis beoogt hiervoor een oplossing te bieden door een tool te ontwikkelen die geografische data analyseert en visualiseert binnen verschillende zoekgebieden.\bigskip

De methodologie voor het bekomen van deze tool bestaat uit vier stappen. Eerst werden via interviews met ET-onderzoekers de benodigde datasets en de vereisten van de onderzoekers geïdentificeerd. Vervolgens werd een vergelijking gemaakt van bestaande GIS-software en bibliotheken, waarbij aandacht werd besteed aan de noden van de ET-onderzoeksgroep. 
QGIS werd hieruit geselecteerd als basis voor de tool vanwege zijn open source karakter, flexibiliteit, ingebouwde Python scripting en PyQGIS, een Python bibliotheek die QGIS en zijn functionaliteit linkt met Python. Hierna werd een proof of concept (PoC) ontwikkeld met een grafische gebruikersinterface (GUI), die gebruikers toelaat om verschillende datasets te analyseren binnen een gedefinieerd of een door de tool gegenereerd zoekgebied. 
De tool biedt opties voor volledige of gedeeltelijke analyses (bv.\ gefilterd op specifieke gewassen). 
De resultaten worden opgeslagen in een in-memory datalaag en kunnen geëxporteerd worden in zowel comma-separated values (CSV) voor verdere analyse als portable document format (PDF) voor rapportage. 
Tenslotte werd een risicoanalyse uitgevoerd om potentiële blind spots, zoals ontbrekende data of technische beperkingen, te identificeren, zodat onderzoekers inzicht krijgen in de beperkingen van de data en op basis daarvan beter geïnformeerde besluiten kunnen trekken\bigskip

De ontwikkelde tool biedt een flexibele en gebruiksvriendelijke oplossing voor het analyseren van geografische data, met aandacht voor modulariteit en uitbreidbaarheid. De GUI is ontworpen om ook gebruikers zonder kennis van GIS te ondersteunen, wat de toepasbaarheid vergroot. Beperkingen zijn dat de tool zich focust op enkel geografische analyse en nog geen berekeningen van energiepotentieel bevat. Bovendien is de kwaliteit van de resultaten afhankelijk van de beschikbare datasets. Mogelijke toekomstige uitbreidingen omvatten de integratie van geavanceerde analysemethoden, zoals energieproductiemodellen, combinatiemethodes voor vergelijkbare data (e.g.\ landbouwkaarten van verschillende instanties) en multi-criteria decision analysis (MCDA).

Deze thesis toont aan dat op GIS-gebaseerde tools een waardevolle bijdrage kunnen leveren aan complexe energieplanningsprocessen, niet alleen voor de ET, maar ook voor andere grootschalige infrastructuurprojecten. Door het combineren van open source technologieën en gebruiksvriendelijke interfaces, biedt de tool een praktische oplossing voor onderzoekers en stakeholders, met ruimte voor verdere ontwikkeling en aanpassing aan nieuwe inzichten.
